\section{Ανάπτυξη  μοντέλων μηχανικής μάθησης}



Σε αυτό το κεφάλαιο, θα αναφερθούμε πρώτα στα μοντέλα που χρησιμοποιήθηκαν,
έπειτα στο πως δημιουργήσαμε το συνολικό pipeline των μοντέλων μηχανική μάθησης
και τέλος στα αποτελέσματα που μας δόθηκαν.

Τον όρος pipeline στην περίπτωση μας τον χρησιμοποιούμε για να εκφράσουμε το σύνολο διεργασιών που πραγματοποιούμε
για να μεταμορφώσουμε τα δεδομένα των \texttt{userInputData} και \texttt{timeSeries} 
σε μορφή δεδομένων που μπορεί να τα δεχτεί η scikit-learn για τα μοντέλα μηχανικής μάθησης που εμπεριέχει,
να εφαρμόσουμε τα μοντέλα με την προσέγγιση που έχουμε. 

\subsection{Προσέγγιση που ακολουθήθηκε και μοντέλα που χρησιμοποιήθηκαν}

Η αρχική υπόθεση που κάναμε για το που θα βασιστούν τα μηχανικά μοντέλα ήταν στην χρονική καθυστέρηση που προέκυπτε από
την απόσταση του αισθητήρα με την πηγή. 
Δηλαδή ότι η αύξηση που θα παρατηρούσαμε μια x χρονική στιγμή για μία πηγή που είναι μισό μέτρο από τον αισθητήρα, θα προέκυπτε 
σε x+n δευτερόλεπτα, κάτι που όπως φάνηκε στην οπτικοποίηση δεν φάνηκε να ισχύει.

Θα δοκιμάσουμε παρόλα αυτά να φέρουμε την πρώτη περίοδο αύξησης που φαίνεται να έχει μεγαλύτερη παράγωγο από προηγούμενες τιμές 
σε ένα χρονικό σημείο που θεωρούμε ότι ιδεατά θα έπρεπε να αυξάνεται τότε.

Από την αρχή επιλέξαμε να χρησιμοποιήσουμε ένα συγκεκριμένα αριθμό δευτερολέπτων από τα πειράματα,
συγκεκριμένα από 30 δευτερόλεπτα πριν την εισαγωγή ρύπου μέχρι 299 δευτερόλεπτα μετά από την εισαγωγή ρύπου.
Όσα πειράματα δεν είχαν τόσα πολλά δεδομένα, είτε πριν είτε μετά, τα αφαιρέσαμε από τα μοντέλα μηχανικών μοντέλων.

Το εύρος δευτερολέπτων  επιλέχθηκε για να έχουμε αιχμαλωτίσει την περίοδο ηρεμίας, του γεγονότος της εισαγωγής ρύπου 
και την μετέπειτα εξέλιξη της τιμής του VOC, χωρίς να αφήσουμε πολύ χρόνο να περάσει.
Επίσης, η μεγάλη πλειοψηφία των πειραμάτων είχε δημιουργηθεί να βρίσκεται τουλάχιστον να εμπεριέχει τουλάχιστον αυτό το εύρος δευτερολέπτων.
Οπότε, είχαμε για κάθε πείραμα 330 δευτερόλεπτα δεδομένων, δηλαδή 5,5 λεπτά, καθώς η περίοδο υποβολής της εισαγωγής πηγής ρύπου θεωρείται ότι 
γίνεται στο μηδενικό δευτερόλεπτο.

Αυτή η επιλογή δυστυχώς ήταν σχετικά αυθαίρετη ,καθώς δεν βρέθηκαν πληροφορίες που να μας λένε πόσο χρόνο παίρνει το αλκοόλ να εξατμιστεί
και τα δεδομένα δεν μας διαφώτισαν πολύ στο ποια περίοδο να κρατήσουμε.
Το εύρος αυτό μας εξασφαλίζει ότι θα είμαστε σίγουροι ότι οι αισθητήρες αντιδρούν στο αλκοόλ που εξατμίζεται στον αέρα
και όχι μόνο το γεγονός της εισαγωγής στο δωμάτιο.
Επίσης, από τον τρόπο που κάναμε τα πειράματα, ξέρουμε ότι η πλειοψηφία των πειραμάτων είναι μες στο εύρος αυτό, οπότε δεν θα χρειαστεί
να αφαιρέσουμε πολλά από τα πειράματα που έχουμε κάνει.

Ο λόγος που αφαιρούμε πειράματα με μικρότερο μέγεθος είναι γιατί η scikit-learn απαιτεί το μέγεθος των στηλών να είναι σταθερό
σε όλα τα δείγματα που δώσουμε σαν είσοδο.
Στην περίπτωση μας, το κάθε πείραμα ξεχωριστά είναι μία γραμμή του πίνακα εισόδου στο μοντέλο και τα δευτερόλεπτα είναι οι στήλες.
Μπορούμε εύκολα να αφαιρέσουμε παραπανήσια δευτερόλεπτα, αλλά δεν είναι εύκολο να δημιουργήσουμε νέα δεδομένα σε χρονικά σημεία 
που δεν έχουμε γνωστές τιμές VOC, καθώς δεν έχουμε κάποια δεδομένη συμπεριφορά.

Για να μπορούσε να δοκιμάσουμε την ιδεατή περίπτωση, μειώσαμε αυτό το εύρος τιμών σε δεκαπέντε δευτερόλεπτα πριν την εισαγωγή πηγής
και 249 δευτερόλεπτα μετά από την εισαγωγή της πηγής,  οπότε 265 δευτερόλεπτα σύνολο.

Επειδή όσο πιο πολλές στήλες δεδομένων έχουμε, τόσο πιο περίπλοκο γίνεται το πρόβλημα και τείνει να γίνει overfitting,
όταν δεν έχουμε αντίστοιχα μεγάλο αριθμό δειγμάτων, πειραμάτων  στην περίπτωση μας, θα επιθυμούσαμε 
να μειώσουμε τον αριθμό αυτό των στηλών, είτε συμπυκνώνοντας την πληροφορία ρύπων σε λιγότερες στήλες, δηλαδή feature extraction,
είτε επιλέγοντας λιγότερες στήλες, με feature selection.
Δυστυχώς, παίρνοντας κάποιες βασικές στατιστικές μετρικές, μέσο και διασπορά, δεν βελτίωσε τα αποτελέσματα, και επειδή μιλάμε για 
Χρονοσειρές ,δεν υπάρχουν δεδομένα που να μην ταιριάζουν με τα υπόλοιπα.
Επειδή η πλειοψηφία των μηχανικών μοντέλων για χρονοσειρές αφορούν προβλέψεις των μελλοντικών τιμών για την χρονοσειρά 
και όχι εύρεση κάποιας τιμής στόχου με βάση τα δεδομένα ,δεν καταφέραμε να βρούμε πολλές πληροφορίες.

Γενικά ,ένα πρόβλημα που είχαμε να αντιμετωπίσουμε ήταν ότι δεν μπορούσαμε απλώς να βάλουμε τις εισόδους όλων των αισθητήρων
στην ίδια είσοδο ,γιατί δεν θα πρόσθεταν παραπάνω δείγματα (samples) στον πίνακα εισόδου (samples matrix ή design matrix),αλλά
πιο πολλές στήλες και χαρακτηριστικά.

Η κύρια προσέγγιση που ακολουθήσαμε είναι Άντι να αναζητήσουμε την θέση της πηγής ρύπου των 
προσπαθώντας να προβλέψουμε άμεσα τα \texttt{x axis} και \texttt{y axis},για κάθε αισθητήρα χρησιμοποιούμε ένα regression μοντέλο
για να προβλέψουμε την ευκλείδεια απόσταση της πηγής από τον αισθητήρα.
Έπειτα, παίρνουμε τις ευκλείδειες αποστάσεις που έχουν προβλέψει και οι 3 αισθητήρες, παίρνουμε τις ακτίνες τους 
και βρίσκομαι σαν θέση πηγής το σημείο που τέμνονται και οι τρεις κύκλοι.
Αυτός ο τρόπος εύρεσης της θέσης πηγή ρύπου λέγεται μέθοδος τριπλευρισμού (trilateration) και χρησιμοποιείται
για την εύρεση θέσης μέσω gps \cite{trilateration}.

Προσεγγίσαμε διάφορα προσεγγίσεις, που όμως δεν αποδείχθηκαν καλύτερες από την παραπάνω μέθοδο τριπλευρισμού, οι οποίες υπάρχουν σε διάφορα αρχεία μες στο project.
Από την πλευρά του regression, δοκιμάσαμε για κάθε αισθητήρα να δημιουργήσουμε ξεχωριστό regression μοντέλο για να προβλέψουμε την \texttt{x axis} και διαφορετικό για την \texttt{y axis}
και να συνδυάσουμε τα αποτελέσματα με το \texttt{StackingRegressor} της \texttt{Scikit\-Learn}.
Η κύρια δοκιμή ήταν με τα δεδομένα να διατηρηθούν όπως είναι κανονικά, αλλά προσπαθήσαμε να μειώσουμε τις στήλες,
ομαδοποιώντας ανά δεκαπέντε δευτερόλεπτα σε ένα κουτί, αντλώντας από αυτό το κουτί την μέση και την διασπορά.
Ακόμα, δοκιμάσαμε με την μέθοδο του τριπλευρισμού να κρατήσουμε από κάθε πείραμα μόνο την μέση και διασπορά, όπως και να αφαιρέσουμε τα πρώτα δευτερόλεπτα ηρεμίας.

Από την πλευρά του classification ,χωρίσαμε τις 16 θέσης πηγής ρύπων σε ξεχωριστές κλάσεις και προσπαθήσαμε να βρούμε
μοντέλο για τον κάθε αισθητήρα ξεχωριστά και μετά ένα συνολικό μοντέλο από τις εξόδους των βασικών μοντέλων με το \texttt{StackingClassifier},
όπως κάναμε και για το \texttt{StackingRegressor}.
Δοκιμάσαμε τον \texttt{VotingClassifier},που αντί να εκπαιδευτεί νέο μοντέλα πάνω στα υπάρχοντα μοντέλα, βγαίνει σαν έξοδο όποια κλάση αποφασίστηκε 
από την πλειοψηφία των βασικών μοντέλων.

Επιπλέον, δοκιμάσαμε να εισάγουμε όλα τα δεδομένα των αισθητήρων σε ένα μοντέλο μόνο, τοποθετώντας πρώτα όλες τις στήλες τους πρώτου αισθητήρα ,μετά του δεύτερου
και τελευταίο του τρίτου.
Πραγματοποιήσαμε classification με κλάσεις που αντιστοιχούσαν στις τέσσερις τιμές του  \texttt{x axis}  και του \texttt{y axis} αντίστοιχα,
διαφορετικά για κάθε αισθητήρα και άξονα και εφαρμόζοντας \texttt{VotingClassifier} για να βρούμε το ιδανικό \texttt{x axis} και \texttt{y axis}.
Τέλος ,αποπειραθήκαμε να δούμε αν μπορεί να εκπαιδεύσουμε το classifier να μαντεύει μεταξύ μακρινής και κοντινής θέσης πηγής, αλλά δεν ήμασταν επιτυχείς σε αυτό.

Τα regression μοντέλα που εξετάστηκαν κάθε φορά για να βρεθεί το κατάλληλο μοντέλο και παράμετροι μέσω hyperparameters tuning, αναφέρονται στο τραπέζι \autoref{tab:models data}.
Αναφέρουμε μόνο τα regression μοντέλα, γιατί την πλειοψηφία αυτών των μοντέλων θα χρησιμοποιήσουμε στο τέλος.
Δυστυχώς, από την βιβλιογραφία δεν βρήκαμε κάποια συγκεκριμένη πρόταση, οπότε κινηθήκαμε σε κάποια δημοφιλή και γνωστά παραδοσιακά μοντέλα μηχανικής μάθησης.

Δεν χρησιμοποιήσαμε μοντέλα βαθιάς μηχανικής μάθησης, γιατί απαιτούν την χρήση ειδικής βιβλιοθήκης για τα νευρωνικά μοντέλα.
Ακόμα, κρίνοντας από τα αποτελέσματα των μοντέλων μας, πιθανότατα δεν θα βοηθούσανε και πολύ, γιατί θα καταλήγανε σε μοντέλα με πολύ μεγάλο variance, λόγω των λίγων δειγμάτων που έχουμε.
Όλα τα μοντέλα μηχανικής μάθησης προσφέρονται από την \texttt{Scikit\-Learn}.

\setlength{\tabcolsep}{4pt}

\begin{table}[h!]
\centering
\caption{Πληροφορίες για τα μοντέλα και τις παραμέτρους που χρησιμοποιήθηκαν}
\label{tab:models_data}

\begin{tabularx}{\textwidth}{|
>{\raggedright\arraybackslash}p{3cm} |
>{\raggedright\arraybackslash\small}p{4cm} |
X |
}
\hline
\textbf{Μοντέλο} & 
\textbf{Hyperparameters} & 
\textbf{Σχόλια} \\
\hline

Linear Regression 
& -- 
& Αποτελεί το πιο βασικό μοντέλο regression, που προσπαθεί να βρει γραμμική λύση για την πρόβλεψη των δεδομένων. Χρησιμοποιείται σαν βάση για τα δύο επόμενα μοντέλα. \\
\hline

Ridge Regularization
&
\makecell[l]{
alpha: [0.01, 0.1, 1,\\ 10, 100] \\
solver: ['auto', 'svd',\\ 'cholesky', 'lsqr']
}
&
Αποτελεί το L2 Regularization του Linear Regression για μείωση του variance. Δεν βελτιώνει ιδιαίτερα το Linear Regression, καθώς οι στήλες δεν συσχετίζονται έντονα μεταξύ τους. \\
\hline

Lasso Regularization
&
\makecell[l]{
max\_iter: [10000] \\
alpha: [0.001, 0.01,\\ 0.1, 1, 10]
}
&
Αποτελεί την L1 Regularization του Linear Regression. Δουλεύει καλύτερα από το Ridge γιατί μπορεί να αφαιρέσει τελείως μη σχετικές στήλες. \\
\hline

ElasticNet Regularization
&
\makecell[l]{
max\_iter: [10000] \\
alpha: [0.001, 0.01,\\ 0.1, 1, 10] \\
l1\_ratio: [0.1, 0.5,\\ 0.7, 0.9]
}
&
Προσπαθεί να βρει ισορροπία μεταξύ L1 και L2 Regularization. Συνήθως αποδίδει χειρότερα σε αυτό το dataset. \\
\hline

Support Vector Regression
&
\makecell[l]{
kernel: ['linear','rbf'] \\
C: [0.1, 1, 10, 100] \\
epsilon: [0.001, 0.01,\\ 0.05, 0.1, 0.2] \\
gamma: ['auto', 1e-3,\\ 1e-2, 1e-1, 1, 'scale']
}
&
Στις περισσότερες περιπτώσεις δίνει τα καλύτερα αποτελέσματα χάρη στις μη γραμμικές λύσεις του rbf kernel. Παρά τους πολλούς hyperparameters είναι σχετικά γρήγορο. \\
\hline

Random Forest Regressor
&
\makecell[l]{
random\_state: [42] \\
n\_estimators: [10, 50,\\ 200, 400] \\
max\_depth: [None, 5,\\ 30, 50] \\
min\_samples\_split: [2, 5] \\
min\_samples\_leaf: [1, \\ 2, 10]
}
&
Αποτελεί γενικά πολύ καλή επιλογή, όμως εδώ δεν ξεπερνά το Support Vector Regressor λόγω πολλών μη χρήσιμων στηλών. \\
\hline

Gradient Boosting Regressor
&
\makecell[l]{
random\_state: [42] \\
n\_estimators: [10, 50,\\ 200, 400] \\
learning\_rate: [0.01, 0.05,\\ 0.1, 0.2] \\
max\_depth: [None, 5,\\ 30, 50] \\
min\_samples\_split: [2, 5] \\
min\_samples\_leaf: [1, \\ 2, 10]
}
&
Προσπαθεί να βελτιώσει το accuracy χτίζοντας διαδοχικά δέντρα. Είναι πιο αργό στο hypertuning και πιο ευαίσθητο σε overfitting σε σχέση με το Random Forest. \\
\hline

\end{tabularx}
\end{table}


Έχουμε τρία σύνολα διεργασιών που θα μπορούσαμε να ονομάσουμε pipeline.
Το πρώτο pipeline είναι το pipeline της \texttt{Scikit\-Learn},
To δεύτερο pipeline,  είναι θα λέγαμε το κύριο ,είναι αυτό που περιγράψαμε στην αρχή του \autoref{sub:machine Models},όπου βρίσκουμε τα καλύτερο μηχανικό μοντέλα για συγκεκριμένες ρυθμίσεις και δεδομένα, περιλαμβάνοντας το pipeline της texttt{Scikit\-Learn}.
Το τρίτος είδος pipeline, που το ονομάζουμε \texttt{loopTroughFunctions},εμπεριέχει το κύριο pipeline και μέσω αυτού, κάνει apply  συναρτήσεις που ονομάζουμε συναρτήσεις φιλτραρίσματος,
που στην περίπτωση μας είναι το πρώτα η διαχείριση των outliers και έπειτα αν θα πάρουμε το gradient ή όχι.
Μετά από αυτό, μας επιστρέφει πίσω τα μετρικές αξιολόγησης που παίρνουμε από το βασικό pipeline και της επιστρέφουμε πίσω σε ένα dataframe που έχει φτιαχτεί για να μας παρουσιάζει σε ένα συγκεκριμένο σημείο τις τελικές μετρικές.

Οι μετρικές εκτίμησης για  τα τελικά αποτελέσματα του συνολικού μοντέλου που έχουμε κατασκευάσει είναι:

\begin{enumerate}
    \item Το R\textsuperscript{2},δηλαδή ο συντελεστής προσδιορισμού (Coefficient of determination),αποτελεί στατιστική μετρική που εκφράζει το ποσοστό της διακύμανσης μιας εξαρτημένης μεταβλητής το οποίο ερμηνεύεται από μια ανεξάρτητη μεταβλητή. Ισοδυναμεί με το τετράγωνο του συντελεστή συσχέτισης \cite{r2}.
    Στην περίπτωση μας ,η ανεξάρτητή τιμή είναι οι θέση πηγή ρύπου που έχουν πραγματικά τα πειράματα που έχουμε δώσει για έλεγχο του συνολικού συστήματος του μηχανικού μοντέλου και οι εξαρτημένες τιμές είναι οι τιμές που πρόβλεψε το συνολικό σύστημα των μηχανικών μοντέλων.
    Αυτή η σχέση μεταξύ πραγματικών και προβλεπόμενων τιμών, είναι οι βάσεις και για τις δύο άλλες μετρικές
    \item Το μέσο τετραγωνικό σφάλμα (Mean Squared Error,MSE),που περιγράφει το μέσο τετραγωνοποιημένο σφάλμα μεταξύ προβλεπόμενων και πραγματικών τιμών.
    \item Η μέση ευκλείδεια απόσταση, που περιγράφει την μέση ευκλείδεια απόσταση που έχουν οι προβλεπόμενες θέσεις πηγής με την πραγματικές θέσεις πηγής.
\end{enumerate}

Η MSE και R\textsuperscript{2} αποτελούν πολύ γνωστές και χρησιμοποιούμενες μετρικές για την αξιολόγηση των μηχανικών regression μοντέλων, ενώ η μέση ευκλείδεια απόσταση,
είναι δικιά μας μετρική που εκφράζει το πως το σφάλμα πρόβλεψης μεταφράζεται σε αποστάσεις στον πραγματικό χώρο του δωματίου.
Το MSE και η μέση ευκλείδεια απόσταση είναι πολύ κοντά μεταξύ τους, αλλά χρησιμοποιούμε το MSE γιατί είναι πιο "χαλαρή" στις βελτιώσεις του συστήματος των μηχανικών μοντέλων,
σε σχέση με την  μέση ευκλείδεια απόσταση, καθώς η τετραγωνοποίηση κάνει τις διαφορές κάτω του πιο μικρές.
Έτσι, μας είναι πιο εύκολο να δούμε γρήγορα τις βελτιώσεις των μηχανικών μοντέλων.

Η R\textsuperscript{2} αποτελεί συμπληρωματική των άλλων δύο στα συνολικά μηχανικά μοντέλα. Είναι όμως πολύ χρήσιμη για τα επιμέρους τελικά μηχανικά μοντέλα,
γιατί σχεδιάζοντας την σχέση μεταξύ πραγματικών και προβλεπόμενων,
μπορούμε να δούμε ποιες προβλέψεις των ευκλείδειων αποστάσεων δεν συμφωνούν με τις πραγματικές τιμές και σε τι τρόπο.
Γενικά, συμφωνούν πάντα μεταξύ τους και οι τρεις μετρικές.
